% Image credits for the photographs and AI illustrations of Book 5 (University, Year 3), Indonesian edition.
% Machine-readable license records live in images/book5/CREDITS.md; keep this
% file in sync with frontmatter/image-credits-book5.tex (one \item per item there).
% Every credit, author, dataset and licence is kept exactly as in the English file.
\chapter*{Kredit Gambar}

Gambar, peta dan grafik dalam buku ini bersifat orisinal; peta dan grafik yang
dibuat dari data yang diterbitkan menyebutkan sumbernya dalam keterangan gambar
atau dalam daftar ini. Foto digunakan di bawah lisensi bebasnya masing-masing,
dengan terima kasih kepada para penulisnya; tautan sumber lengkap dan URL lisensi
untuk setiap foto tercantum dalam \texttt{images/book5/CREDITS.md} di repositori buku.

\bigskip
Selain foto dan gambar orisinal, beberapa figur adalah ilustrasi yang dihasilkan AI,
dibuat dengan pembuatan gambar OpenAI dari prompt yang ditulis oleh editor buku, dan
ditinjau keakuratan ilmiahnya sebelum dimasukkan. Gambar-gambar ini hanya menampilkan
adegan umum, tidak pernah tempat nyata, orang nyata, atau spesimen nyata. Prompt
lengkap untuk setiap gambar yang dihasilkan tercatat dalam
\texttt{images/book5/ai/PROMPTS.md} di repositori.

\bigskip
\noindent Foto dan data, menurut bab:
\begin{itemize}\small
  \item Bab~1 (seismologi dan pencitraan Bumi): PREM, A.~M.~Dziewonski dan D.~L.~Anderson (1981), sebagaimana ditabulasikan dalam ObsPy; ak135, B.~L.~N.~Kennett, E.~R.~Engdahl dan R.~Buland (1995).
  \item Bab~3 (gravitasi dan bentuk Bumi): satelit kembar GRACE Follow-On, gambaran seniman, NASA/JPL-Caltech, domain publik (Wikimedia Commons). Neraca massa lembaran es dari IMBIE Team (2021), UK Polar Data Centre, Open Government Licence v3.0.
  \item Bab~4 (geomagnetisme dan paleomagnetisme): International Geomagnetic Reference Field, generasi keempat belas, IAGA Division~V Working Group V-MOD, koefisien didistribusikan oleh NOAA NCEI.
  \item Bab~5 (model termal litosfer): umur dasar samudra, M.~Seton dkk. (2020), EarthByte, CC BY 4.0; relief global ETOPO1, NOAA NCEI, domain publik.
  \item Bab~6 (konveksi mantel dan gaya penggerak lempeng): pantai terangkat Ballyhillin, Malin Head, A.~F.~Borchert, CC BY-SA 4.0 (Wikimedia Commons); PREM, A.~M.~Dziewonski dan D.~L.~Anderson (1981).
  \item Bab~7 (mekanisme deformasi dan zona geser): milonit augen, Mikenorton, CC BY-SA 4.0 (Wikimedia Commons).
  \item Bab~9 (tektonik ekstensi dan sesar geser-mendatar): topografi basin-and-range di selatan Death Valley, Ken Lund, CC BY-SA 2.0 (Wikimedia Commons).
  \item Bab~14 (geokimia isotop stabil): tumpukan bentik $\delta^{18}$O LR04, L.~E.~Lisiecki dan M.~E.~Raymo (2005).
  \item Bab~15 (interaksi air--batuan): teras travertin Palette Spring, Mammoth Hot Springs, Jim Peaco, National Park Service, domain publik (Wikimedia Commons). Konstanta kesetimbangan dari basis data PHREEQC phreeqc.dat, D.~L.~Parkhurst dan C.~A.~J.~Appelo, U.S. Geological Survey, domain publik.
  \item Bab~16 (stratigrafi sekuen): Book Cliffs di atas Helper, Utah, David Mackie, domain publik (Wikimedia Commons); tumpukan bentik $\delta^{18}$O LR04, L.~E.~Lisiecki dan M.~E.~Raymo (2005).
  \item Bab~17 (model fasies): stratifikasi silang hummocky di Formasi Arisaig, Michael C.~Rygel, CC BY-SA 3.0 (Wikimedia Commons), dipotong.
  \item Bab~18 (geomorfologi kuantitatif): Ngarai Niagara dan Air Terjun Horseshoe, Ken Lund, CC BY-SA 2.0 (Wikimedia Commons), dipotong.
  \item Bab~19 (dinamika fluida geofisika): angin rata-rata zonal NCEP--NCAR Reanalysis~1, data disediakan oleh NOAA PSL, Boulder, Colorado, AS; World Ocean Atlas 2023, NOAA NCEI.
  \item Bab~20 (dinamika atmosfer): sistem tekanan rendah di lepas pantai Islandia, NASA/GSFC, MODIS Rapid Response Team, Jacques Descloitres, domain publik (Wikimedia Commons); angin rata-rata zonal NCEP--NCAR Reanalysis~1, data disediakan oleh NOAA PSL, Boulder, Colorado, AS.
  \item Bab~21 (dinamika samudra): pelampung Argo yang dilepas dari kapal riset, NOAA Global Ocean Monitoring and Observing Program, domain publik (Wikimedia Commons).
  \item Bab~22 (model iklim): superkomputer NASA Center for Climate Simulation, Goddard Space Flight Center, 2010, domain publik (Wikimedia Commons); suhu udara permukaan dan radiasi puncak atmosfer NCEP--NCAR Reanalysis~1, data disediakan oleh NOAA PSL, Boulder, Colorado, AS.
  \item Bab~24 (glasiologi): gletser keluaran di Greenland barat daya, NASA / Jim Yungel, Operation IceBridge, 2014, domain publik (Wikimedia Commons); elevasi permukaan Greenland dari ETOPO1 Ice Surface, NOAA NCEI; luas es laut Arktik dari NSIDC Sea Ice Index, versi~4 (Fetterer dkk., 2025, doi:10.7265/a98x-0f50).
  \item Bab~25 (penginderaan jauh dan SIG): interferogram 1995--2001 pusat gunung api Three Sisters, Oregon, U.S. Geological Survey (C.~Wicks dkk., 2002), domain publik (Wikimedia Commons); spektrum laboratorium dari ECOSTRESS Spectral Library, versi~1.0 (S.~K.~Meerdink, S.~J.~Hook, D.~A.~Roberts dan E.~A.~Abbott, 2019; A.~M.~Baldridge dkk., 2009), direproduksi atas kebaikan Jet Propulsion Laboratory, California Institute of Technology, Pasadena, California.
  \item Bab~26 (interior dan permukaan planet): Kawah Meteor Barringer, Arizona, USGS/D.~Roddy, domain publik (Wikimedia Commons); Olympus Mons, mosaik pengorbit Viking, NASA/JPL/USGS, domain publik (Wikimedia Commons). Fungsi produksi dan kronologi kawah sebagaimana diberikan dalam perangkat lunak Craterstats karya G.~G.~Michael, mengikuti G.~Neukum, B.~A.~Ivanov dan W.~K.~Hartmann (2001).
  \item Bab~27 (atmosfer planet dan kelayakhunian): jaringan lembah kering di Mars, Mars Express HRSC, ESA/DLR/FU Berlin, CC BY-SA 3.0 IGO (Wikimedia Commons); Ligeia Mare, Titan, mosaik radar Cassini, NASA/JPL/USGS, domain publik (Wikimedia Commons). Pencocokan zona layak huni R.~K.~Kopparapu dkk. (2013); tekanan uap metana dari NIST Chemistry WebBook.
  \item Bab~28 (eksoplanet): HIP~65426~b yang dicitrakan oleh James Webb Space Telescope, NASA/ESA/CSA, A~Carter (UCSC), tim ERS~1386, dan A.~Pagan (STScI), CC BY~4.0 (Wikimedia Commons). Planet terkonfirmasi dari NASA Exoplanet Archive, tabel Planetary Systems Composite Parameters, diakses 23 Oktober 2026: penelitian ini memanfaatkan NASA Exoplanet Archive, yang dioperasikan oleh California Institute of Technology, berdasarkan kontrak dengan National Aeronautics and Space Administration di bawah Exoplanet Exploration Program. Kurva model massa--jari-jari L.~Zeng dkk. (2016, 2019).
  \item Bab~29 (Bima Sakti dan galaksi-galaksi): panorama Bima Sakti proyek GigaGalaxy Zoom, ESO/S.~Brunier, CC BY~4.0 (Wikimedia Commons); galaksi spiral M101, European Space Agency dan NASA (Hubble, heic0602; data Hubble K.~D.~Kuntz, F.~Bresolin, J.~Trauger, J.~Mould dan Y.-H.~Chu; pengolahan citra Davide De~Martin), CC BY~4.0 (Wikimedia Commons). Kecepatan melingkar Galaksi dari A.-C.~Eilers, D.~W.~Hogg, H.-W.~Rix dan M.~K.~Ness (2019); model massa E.~Pouliasis, P.~Di~Matteo dan M.~Haywood (2017); Sefeid klasik OGLE di Awan Magellan Besar, I.~Soszy\'nski dkk. (2015), Acta Astronomica 65, 297, didistribusikan oleh CDS, Strasbourg (VizieR J/AcA/65/297).
  \item Bab~30 (dari mana unsur-unsur berasal): Nebula Kepiting, NASA, ESA, J.~Hester dan A.~Loll (Arizona State University), domain publik (Wikimedia Commons); kilonova GW170817 di NGC~4993, NASA dan ESA, ucapan terima kasih kepada A.~J.~Levan, N.~R.~Tanvir, A.~Fruchter dan O.~Fox (Hubble, heic1717a), CC BY~4.0 (Wikimedia Commons). Kelimpahan Matahari M.~Asplund, A.~M.~Amarsi dan N.~Grevesse (2021); penampang tangkapan neutron dari KADoNiS v0.3 (I.~Dillmann, R.~Plag, F.~K\"appeler dan T.~Rauscher); waktu paruh dari IAEA Live Chart of Nuclides (ENSDF); metalisitas survei Geneva--Copenhagen sebagaimana dianalisis ulang oleh L.~Casagrande dkk. (2011), didistribusikan oleh CDS, Strasbourg (VizieR J/A+A/530/A138).
\end{itemize}
